\documentclass[12pt,letterpaper]{article}   % Tamaño 12, papel carta

% Paquetes esenciales
\usepackage[spanish]{babel}        % Idioma español
\usepackage[utf8]{inputenc}        % Codificación UTF-8
\usepackage[T1]{fontenc}           % Codificación de fuente
\usepackage{setspace}              % Control del interlineado
\usepackage[margin=2.5cm]{geometry}% Márgenes de 2.5cm
\usepackage{parskip}               % Evita sangrías y añade espacio entre párrafos

\usepackage{graphicx}              % Para incluir imágenes 
\usepackage{titling}               % Para personalizar la portada
\usepackage{amsmath}               % Para mates
\usepackage{amssymb}               % Letras de mates
\usepackage{hyperref}              % Para el índice

\usepackage{amsfonts}
\usepackage{pgfplots}
\usepackage{booktabs}

% Interlineado 1.5
\onehalfspacing

% Justificación del texto
\usepackage{ragged2e}
\justifying

% No sangría en párrafos
\setlength{\parindent}{0pt}

% Numerado de pagina arriba a la derecha
\usepackage{fancyhdr}
\pagestyle{fancy}
\fancyhf{} % Limpia encabezados y pies predeterminados
\fancyhead[R]{\thepage} % Número de página en encabezado derecha
\renewcommand{\headrulewidth}{0pt} % Elimina la línea bajo el encabezado

% Que los titulos de las secciónes tengan el mismo tamaño
\usepackage{titlesec}

% Configuración de \section
\titlespacing*{\section}{0pt}{2ex plus 1ex minus .2ex}{-0.5ex} 
\titleformat{\section}[block]{\normalfont\normalsize\bfseries}{\thesection}{1em}{}

\titlespacing*{\subsection}{0pt}{2ex plus 1ex minus .2ex}{-0.5ex} 
\titleformat{\subsection}[block]{\normalfont\normalsize\bfseries}{\thesubsection}{1em}{}

\titlespacing*{\subsubsection}{0pt}{2ex plus 1ex minus .2ex}{-0.5ex} 
\titleformat{\subsubsection}[block]{\normalfont\normalsize\bfseries}{\thesubsubsection}{1em}{}

\titlespacing*{\paragraph}{0pt}{2ex plus 1ex minus .2ex}{-0.5ex} 
\titleformat{\paragraph}[block]{\normalfont\normalsize\bfseries}{\paragraph}{1em}{}

\titlespacing*{\subparagraph}{0pt}{2ex plus 1ex minus .2ex}{-0.5ex} 
\titleformat{\subparagraph}[block]{\normalfont\normalsize\bfseries}{\subparagraph}{1em}{}


% Consola
\usepackage{listings}
\usepackage{xcolor}
\definecolor{consolegray}{gray}{0.95}
\definecolor{consoleborder}{gray}{0.7}
\lstnewenvironment{smallconsole}[1][]{
  \lstset{
    basicstyle=\ttfamily\tiny, 
    backgroundcolor=\color{consolegray},
    frame=single,
    rulecolor=\color{consoleborder},
    breaklines=true,
    columns=fullflexible,
    showstringspaces=false,
    captionpos=b,
    xleftmargin=1em,
    xrightmargin=1em,
    aboveskip=1em,
    belowskip=1em,
    #1 % ← Para pasar opciones como caption
  }
}{}


% Tamaño de las captions
\usepackage{caption}
\captionsetup{font=scriptsize}

% Multicolumnas
\usepackage{multicol}
\usepackage{listings}
\usepackage{xcolor}
\usepackage{multicol}

\definecolor{consolegray}{gray}{0.95}
\definecolor{consoleborder}{gray}{0.7}

% Define el estilo una vez
\lstdefinestyle{smallconsolestyle}{
  basicstyle=\ttfamily\tiny,
  breaklines=true,
  columns=fullflexible,
  showstringspaces=false,
  captionpos=b,
  xleftmargin=0em,
  xrightmargin=0em,
  aboveskip=1em,
  belowskip=1em,
  frame=none,
  backgroundcolor=\color{white},
  rulecolor=\color{white}
}




% Configuración para usar el punto como separador decimal
\renewcommand{\theequation}{\arabic{equation}}



\begin{document}




% ----------------------------------
% Portada
% ----------------------------------
\begin{titlepage}
    \begin{center}
        \vspace*{2cm}
        \textbf{Portafolio para Criptomonedas}\\[2cm]

        Heriberto Espino Montelongo - 175199\\[0.5cm]
        Universidad de las Américas Puebla\\[0.5cm]
        P25-LEC4032-2: Econometría II\\[0.5cm]
        Dra. Daniela Cortés Toto\\[0.5cm]
        30 de abril de 2025\\
    \end{center}
    \vfill
\end{titlepage}


% ----------------------------------
% Contenido del documento
% ----------------------------------


\pagenumbering{Roman}
\setcounter{page}{1}
\tableofcontents
\newpage



% A partir de aquí empieza la numeración arábiga desde 1
\pagenumbering{arabic}
\setcounter{page}{1}




% \begin{figure}[h!]
%     \centering
%     \includegraphics[width=0.7\textwidth]{imagenes/01-01-precipitacion.pdf}
%     \caption{\textit{Elaboración Propia}. Descripción de la imagen.}
%     \label{fig:mi_imagen}
% \end{figure}









\section{Introducción}
\subsection{Descripción del contexto y la importancia del problema seleccionado}
En los últimos años, el ecosistema de criptomonedas ha evolucionado de manera acelerada, alcanzando un valor de mercado superior a los 2,billones de dólares a principios de 2025. Esta dinámica ha atraído tanto a inversionistas institucionales como a usuarios minoristas, pero ha puesto de manifiesto la elevada volatilidad y correlación existente entre los distintos activos digitales. El análisis adecuado de portafolios de criptomonedas requiere metodologías sólidas que permitan evaluar el compromiso riesgo-retorno, así como herramientas visuales intuitivas que faciliten la toma de decisiones basadas en datos.\newline

\subsection{Delimitación del problema, en términos temporales, espaciales y conceptuales}
Este trabajo se centra en el periodo comprendido entre mayo de 2024 y mayo de 2025, y considera un conjunto diverso de criptomonedas representativas seleccionadas por criterios de capitalización de mercado, funcionalidad y relevancia en el ecosistema:

\begin{enumerate}
\item \textbf{Las más grandes por capitalización de mercado}: BTC, ETH, USDT, BNB, USDC, XRP, ADA, SOL, DOGE.
\item \textbf{Stablecoins relevantes}: USDT, USDC, DAI, USDS.
\item \textbf{Tokens de plataformas blockchain}: TRX, TON, MATIC, DOT, AVAX, STETH, WSTETH, SUI.
\item \textbf{Activos envueltos (wrapped assets)} para interoperabilidad: WBTC, WTRX, WSTETH.
\item \textbf{Criptos con adopción en pagos o DeFi}: LTC, LINK, XLM, HBAR, BCH.
\item \textbf{Memecoins populares}: DOGE, SHIB.
\item \textbf{Tokens de exchanges}: LEO.
\end{enumerate}

Se analizan dos estrategias de inversión: compras directas (spot) y contratos por diferencia (CFD), con el objetivo de comparar su comportamiento, eficiencia y riesgo bajo condiciones de mercado reales y simuladas. Geográficamente, las operaciones se asumen realizadas en exchanges globales con datos de precios y volúmenes asociados a transacciones en dólares estadounidenses. Conceptualmente, el estudio abarca técnicas de selección de portafolios (frontera eficiente de Markowitz), modelado estocástico (Movimiento Browniano Geométrico) y análisis de correlaciones avanzadas.

\subsection{Descripción de los objetivos del proyecto}
Los objetivos específicos de este proyecto son:
\begin{itemize}
\item Construir y comparar dos portafolios de criptomonedas (spot y CFD) mediante la frontera eficiente de Markowitz.
\item Evaluar el comportamiento histórico de los portafolios y proyectar su evolución a corto plazo usando un modelo GBM.
\item Realizar análisis de correlación y clustering para identificar agrupaciones de activos y diversificación potencial.
\item Integrar un análisis de precio y volumen de transacción para BTC, ETH, DOGE y XRP, relacionando eventos del ecosistema (tecnológicos, regulatorios y narrativas de mercado) con movimientos de mercado.
\item Desarrollar un dashboard visual aplicando principios de teoría del color, leyes de Gestalt y percepción visual, para facilitar la interpretación de resultados.
\end{itemize}

\section{Metodolog'{\i}a}
\subsection{Explicación precisa de las herramientas y t'{e}cnicas de ciencia de datos utilizadas}
Se emplearon las siguientes herramientas:
\begin{itemize}
\item \textbf{Lenguaje Python} con librerías \texttt{pandas} para manejo de datos, \texttt{numpy} para cálculos numéricos y \texttt{matplotlib} y \texttt{seaborn} para visualización.
\item \textbf{Optimización de portafolios}: implementación de la frontera eficiente de Markowitz mediante simulaciones Monte Carlo y resolución de programas cuadráticos con \texttt{cvxopt}.
\item \textbf{Modelo estocástico}: Movimientos Brownianos Geométricos (GBM) para simular trayectorias de precios a 30 días, obteniendo distribuciones de posibles precios futuros.
\item \textbf{Análisis de correlación}: cálculo de coeficientes de Pearson sobre log-retornos diarios; derivación de pseudo-métrica de disimilitud $d_{ij}=1-|\rho_{ij}|$.
\item \textbf{Clustering y visualización de grafos}: generación de un clustermap jerárquico y grafo con layout de Kamada-Kawai usando \texttt{scikit-learn} y \texttt{networkx}.
\item \textbf{Análisis de eventos}: correlacionar picos de precio y volumen con anuncios de figuras clave (p.ej., Elon Musk), actualizaciones de protocolo (Ethereum 2.0) y decisiones regulatorias (caso XRP vs SEC).
\item \textbf{Diseño de dashboard}: aplicación de paletas Tab10, Coolwarm y divergentes; organización en formato 16:9, resolución 4K; principios de proximidad, similitud y continuidad de la Gestalt.
\end{itemize}

\subsection{Descripci'{o}n sucinta de los datos utilizados y el proceso de recolecci'{o}n, limpieza y an'{a}lisis aplicado}
Los datos de precios y volúmenes diarios se obtuvieron de APIs públicas de exchanges (CoinGecko, Binance) y bases de datos de calidad financiera. El periodo considerado abarca de mayo de 2024 a mayo de 2025. El proceso incluyó:
\begin{enumerate}
\item \textbf{Recolecci'{o}n}: descarga de series de precios de cierre y volúmenes diarios para cada activo.
\item \textbf{Limpieza}: manejo de valores faltantes por interpolación lineal y verificación de estacionariedad con la prueba de Dickey-Fuller.
\item \textbf{Transformaci'{o}n}: cálculo de log-retornos diarios para asegurar aditividad y log-normalidad.
\item \textbf{An'{a}lisis exploratorio}: estadísticas descriptivas, identificación de outliers y análisis de correlación preliminar.
\end{enumerate}

\section{Resultados}
\subsection{Presentaci'{o}n de los resultados obtenidos mediante herramientas de ciencia de datos}
Se construyeron dos portafolios eficientes, cuyos pesos óptimos se presentan en el Gráfico de Barras (Figura \ref{fig:pesos}). El portafolio spot mostró una asignación diversificada, mientras que el de CFDs concentró mayor peso en activos volátiles como DOGE. En el scatterplot de riesgo (desviación \textit{vs.} retorno), ambos portafolios superaron al benchmark de capitalización de mercado, con el CFD exhibiendo mayor retorno esperado a costa de un riesgo significativamente superior.

La simulación GBM a 30 días permite observar distribuciones de precios futuros (Figura \ref{fig:gbm_hist}). El portafolio CFD presenta una cola inferior pronunciada (posible pérdida total) y una cola superior extendida (multiplicación por hasta 6 veces), en contraste con el spot más conservador (colocándose entre 0.5 y 1.5 veces la inversión).

El análisis de correlaciones reveló una alta relación positiva entre la mayoría de las criptomonedas (clustermap en Figura \ref{fig:clustermap}). TON, DAI y LEO se destacaron como outliers de diversificación, debido a su menor correlación con el grupo principal. El grafo de Kamada-Kawai refuerza esta separación, mostrando distancias métricas que facilitan la interpretación de similitud estructural.

El estudio de precio y volumen para BTC, ETH, DOGE y XRP demostró:
\begin{itemize}
\item \textbf{BTC}: aumento de 60,000 a 100,000 USD y récord de volumen de $1\times10^{11}$, impulsado por adopción institucional.
\item \textbf{DOGE}: subió de 0.10 a 0.45 USD con picos de $4\times10^{10}$ tras declaraciones de Musk.
\item \textbf{ETH}: crecimiento de 0.5 a 3.0 USD asociado a Ethereum 2.0, volumen máximo de $8\times10^{10}$.
\item \textbf{XRP}: osciló de 0.2 a 1.4 USD con volúmenes de $5\times10^{10}$ vinculados a avances legales con la SEC.
\end{itemize}

\subsection{Utilizar gr'{a}ficas o paneles para ilustrar claramente los hallazgos}
% Aquí se insertarían referencias a las figuras del dashboard generado: barras de pesos, líneas de evolución, scatterplot, histograma GBM, clustermap, boxplots y grafo.

\section{Conclusi'{o}n y recomendaciones}
\subsection{Sintetiza la conclusi'{o}n obtenida del proyecto}
El análisis conjunto demuestra que, aunque los CFDs pueden generar mayores retornos esperados, su perfil de riesgo es considerablemente más elevado y puede llevar a pérdidas totales en horizontes cortos. El portafolio spot, en cambio, ofrece estabilidad relativa y rendimientos moderados, siendo más adecuado para inversores aversos al riesgo.

\subsection{Proporciona recomendaciones basadas en los resultados obtenidos}
Se recomienda:
\begin{itemize}
\item Para estrategias de mediano plazo (>6 meses), priorizar un portafolio spot diversificado, incorporando activos con baja correlación (p.ej., TON, DAI).
\item Implementar monitoreo continuo de correlaciones y volúmenes, especialmente ante eventos de alto impacto (anuncios regulatorios, actualizaciones de protocolo).
\item Considerar posiciones en CFDs solo para operadores con alta tolerancia al riesgo y capacidad de absorción de pérdidas.
\end{itemize}

\subsection{Identifica los l'{i}mites del proyecto y sugiere posibles direcciones para mejorarlo}
Las limitaciones incluyen asunciones de normalidad implícita en el GBM, volatilidad constante y dependencia de datos históricos para proyecciones. Futuras investigaciones podrían:
\begin{itemize}
\item Incorporar modelos de volatilidad estocástica (Heston) o saltos de precio (Merton).
\item Ampliar el universo de activos con tokens DeFi, NFT y stablecoins algorítmicas.
\item Desarrollar interfaces interactivas empleando dashboards web (Dash, Streamlit) para análisis en tiempo real.
\end{itemize}

\end{document}